% Source: https://artofproblemsolving.com/community/c3h2221705_system_of_equations
% Stable identity: aops-topic-2221705-q01
\textbf{(a)} The system $x+y+z=3$, $x-y=0$ has the solutions $(t,t,3-2t)$ for every real $t$, for example $(1,1,1)$. The system $x+y=2$, $x-y=0$, $2x+3y=5$ has the solution $(1,1)$. (Also, $n$ equations in $n$ unknowns need not have a solution at all: $x+y=1$, $x+y=2$.) What is true is that a \emph{unique} solution needs at least $n$ independent equations; (b) proves this for $n=3$ and two equations.

\textbf{(b)} Suppose $(x_0,y_0,z_0)$ is a solution. It is enough to find $(p,q,r)\ne(0,0,0)$ with
\[a_1p+b_1q+c_1r=0,\qquad a_2p+b_2q+c_2r=0,\tag{$*$}\]
because then $(x_0+tp,\,y_0+tq,\,z_0+tr)$ is a solution for every real $t$, and these are infinitely many different solutions.

Take $p=b_1c_2-b_2c_1$, $q=c_1a_2-c_2a_1$, $r=a_1b_2-a_2b_1$. Expanding,
\[a_1p+b_1q+c_1r=a_1b_1c_2-a_1b_2c_1+b_1c_1a_2-b_1c_2a_1+c_1a_1b_2-c_1a_2b_1=0,\]
and in the same way $a_2p+b_2q+c_2r=0$. If $(p,q,r)\ne(0,0,0)$ we are done.

Otherwise $p=q=r=0$. If $(a_1,b_1,c_1)=(0,0,0)$, then $(*)$ is the single equation $a_2p+b_2q+c_2r=0$. If not, say $a_1\ne0$ (the cases $b_1\ne0$, $c_1\ne0$ are the same with the letters permuted); then $r=0$ and $q=0$ give $b_2=\dfrac{a_2}{a_1}b_1$ and $c_2=\dfrac{a_2}{a_1}c_1$, so the second equation of $(*)$ is $\dfrac{a_2}{a_1}$ times the first, and $(*)$ is again a single equation. A single equation $\alpha p+\beta q+\gamma r=0$ always has a nonzero solution: $(1,0,0)$ if $\alpha=\beta=\gamma=0$; otherwise $(\beta,-\alpha,0)$ if $\alpha\ne0$ or $\beta\ne0$, and $(\gamma,0,-\alpha)=(\gamma,0,0)$ if $\alpha=\beta=0\ne\gamma$.

So in every case a solution comes with infinitely many others, and the system never has exactly one solution. $\blacksquare$
